\section{Find Pairs with a Given Sum in a Sorted DLL}
\label{sec:ll-pairs-sum-dll}

\problemstatement{Given a doubly linked list sorted in ascending order and a
target, find all pairs of nodes whose values sum to the target, using $O(1)$
extra space.}

\begin{patternbox}
A sorted DLL supports the array \textbf{two-pointer} technique directly: one
pointer at the head (smallest), one at the tail (largest). The
\texttt{prev} pointer lets the right pointer move backward, which a singly
linked list cannot do.
\end{patternbox}

\subsection*{Converge from both ends}

Place \texttt{left} at the head and \texttt{right} at the tail. Compare
their sum to the target: if equal, record the pair and move both inward; if
too small, advance \texttt{left} (via \texttt{next}) to increase the sum; if
too large, retreat \texttt{right} (via \texttt{prev}) to decrease it. Stop
when the pointers meet or cross. Sortedness guarantees this monotonic
convergence finds every pair without revisiting.

\begin{ideabox}[The prev Pointer Enables Backward Two Pointers]
The classic two-sum-on-sorted-array needs a right pointer that moves left.
A DLL's \texttt{prev} provides exactly that, so the linear two-pointer sweep
transfers unchanged from arrays to sorted doubly linked lists.
\end{ideabox}

\subsection*{Algorithm}

\begin{enumerate}
  \item \texttt{left = head}; \texttt{right =} last node.
  \item While \texttt{left != right} and not crossed: compare
        \texttt{left->val + right->val} to the target; record on equality
        and move both, else move the appropriate pointer.
  \item Output all recorded pairs.
\end{enumerate}

\subsection*{Dry run}

\dryrun{Sorted DLL $1,2,3,4,6$; target $6$.}

\begin{itemize}
  \item $(1,6)$? $7>6$, move right to $4$. $(1,4)$? $5<6$, move left to
        $2$. $(2,4)=6$ -- record.
  \item Move both; pointers cross. Pair found: $(2,4)$.
\end{itemize}

\subsection*{C++ implementation}

\begin{lstlisting}
#include <bits/stdc++.h>
using namespace std;

struct Node {
    int val; Node* prev; Node* next;
    Node(int v) : val(v), prev(nullptr), next(nullptr) {}
};

Node* build(const vector<int>& a) {
    Node dummy(0), *t = &dummy;
    for (int x : a) { Node* n = new Node(x); t->next = n; n->prev = t; t = n; }
    Node* head = dummy.next; if (head) head->prev = nullptr;
    return head;
}

vector<pair<int,int>> pairsWithSum(Node* head, int target) {
    vector<pair<int,int>> res;
    if (!head) return res;
    Node* left = head;
    Node* right = head;
    while (right->next) right = right->next;       // go to tail
    while (left != right && right->next != left) {
        int sum = left->val + right->val;
        if (sum == target) {
            res.push_back({left->val, right->val});
            left = left->next; right = right->prev;
        } else if (sum < target) left = left->next;
        else right = right->prev;
    }
    return res;
}

int main() {
    for (auto& [a, b] : pairsWithSum(build({1, 2, 3, 4, 6}), 6))
        cout << "(" << a << "," << b << ") ";
    cout << "\n";                                  // (2,4)
    return 0;
}
\end{lstlisting}

\begin{complexitybox}
\textbf{Time Complexity.} $O(n)$ after reaching the tail (also $O(n)$).
\textbf{Space Complexity.} $O(1)$ beyond the output.
\end{complexitybox}

\begin{takeawaybox}
Sorted-DLL pair-sum is the two-pointer two-sum with the right pointer moving
backward via \texttt{prev}. The DLL's backward link is precisely what lets
this array technique work on a list.
\end{takeawaybox}
